\section*{ID6595: Kissing-Geometry IFS on the Sphere (Algorithmic Description): R\textasciicircum{}d}
\paragraph{Metadata.}
PiID: 2701021835646 | RTEID: RTE:P25609.8496:T4.4980 | UUID: 40ed4f71-d4fd-5b65-bdf8-0d8b0d4f37b7

\paragraph{Canonical Statement.}
1. Start with a unit sphere \textbackslash{}(S\textasciicircum{}2\textbackslash{}) (or unit ball in \textbackslash{}(\textbackslash{}mathbb\{R\}\textasciicircum{}d\textbackslash{})), generation \textbackslash{}(k=0\textbackslash{}) single central object of radius \textbackslash{}(R\_0\textbackslash{}). 2. For each object at generation \textbackslash{}(k\textbackslash{}) of radius \textbackslash{}(R\_k\textbackslash{}) place \textbackslash{}(N\textbackslash{}) child objects of radius \textbackslash{}(R\_\{k+1\}=r R\_k\textbackslash{}) centered at the canonical **kissing positions** relative to the parent (i.e., the points where \textbackslash{}(N\textbackslash{}) equal spheres touch the parent). On \textbackslash{}(S\textasciicircum{}2\textbackslash{}) these are the \textbackslash{}(N\textbackslash{}) unit vectors forming the kissing configuration. 3. For each child, repeat: place its own \textbackslash{}(N\textbackslash{}) children scaled by \textbackslash{}(r\textbackslash{}), centered at the child’s local kissing geometry (you can rotate the local kissing template to align it to the child’s tangent frame). 4. Iterate infinitely. The TZP fractal \textbackslash{}(\textbackslash{}mathcal\{F\}\textbackslash{}) is the limit set of centers (or union of spheres, or union of scaled patches) under infinite recursion.

\paragraph{Canonical Equation.}
\[
\mathbb{R}^d
\]

\paragraph{Dictionary.}
Kissing-Geometry IFS on the Sphere (Algorithmic Description): R\textasciicircum{}d — R\textasciicircum{}d

\paragraph{Encyclopedia.}
Kissing-Geometry IFS on the Sphere (Algorithmic Description): R\textasciicircum{}d is a symbolic expression, written R\textasciicircum{}d. It introduces a symbol or relation used elsewhere in the framework. Within the Trawin Zero Point registry it serves as one element of the framework's mathematical narrative.
